\chapter{Introdución}
\label{chap:introducion}

\section{Contexto e motivación}
\label{sec:contexto}


\lettrine{A}{} familia de consolas Game Boy é unha das máis populares e exitosas da historia dos videoxogos. Ata 2020 seguiuse producindo hardware capaz de executar de forma nativa xogos da familia Advance, a Nintendo 3DS.
O seu percorrido e éxito é polo tanto longo e a súa influencia na industria dos videoxogos e na cultura popular é indiscutible.

\vspace{1em}
A pesar da fin do seu ciclo de vida comercial, esta consola e os seus xogos seguen sendo relevantes para unha grande comunidade de afeccionados e desenvolvedores.
Por exemplo, séguense a lanzar xogos non oficiais comerciais como o JETT RIDER Mini Hero, un xogo para PC e Game Boy que conseguiu un enorme éxito na súa campaña de kickstarter,
que quintiplicou a súa meta inicial de financiamento % \cite{jettriderminihero}.

\vspace{1em}
Para as demandas da situación actual do desenvolvemento de videoxogos para esta plataforma e para a preservación da súa herdanza cultural a ferramenta principal é a emulación.
A emulación é unha técnica que permite a execución de software deseñado para unha plataforma nunha plataforma diferente, mediante a creación dun software que simula o comportamento do hardware orixinal.
No caso da Game Boy a emulación é completa, xa que a emulación non se limita a realizar unha emulación dunha parte do hardware, senón que todos os compoñentes do sistema son emulados.
Existen unha gran cantidade de emuladores de Game Boy, para diferentes plataformas e con diferentes características, tanto propietarios como libres e de código aberto.
Todos estes emuladores emplegan como base documental a gran cantidade de documentación producida pola comunidade a través da análise do hardware mediante enxeñería inversa.

\vspace{1em}
Grazas a este traballo de documentación e enxeñería inversa da comunidade,
o desenvolvemento dun emulador de Game Boy é unha tarefa abordable que combina retos de baixo nivel
con conceptos transversais da arquitectura de computadores.
Esta combinación convérteo nun exercicio especialmente axeitado tanto para a aprendizaxe como para a experimentación con sistemas reais.

\vspace{1em}
Neste contexto, o presente traballo formúlase cun dobre obxectivo.
Por unha banda, perségue unha finalidade didáctica:
servir como medio para afondar de xeito práctico no funcionamento interno da consola e nos principios da emulación,
documentando as decisións de deseño tomadas ao longo do desenvolvemento.
Por outra banda, propón unha finalidade técnica: ofrecer unha solución en forma de librería que poida ser integrada de forma doada en distintos contextos,
con especial atención á súa execución en sistemas embebidos de recursos limitados.
A partir dese núcleo, o traballo materialízase nunha consola portátil de referencia
construída sobre hardware accesible e de baixo custo e documentada,
concibida como unha alternativa á escaseza do hardware orixinal
e ás solucións comerciais caras e/ou pechadas.
Este posicionamento desenvólvese en detalle na sección~\ref{sec:posicionamento}.

\section{Obxectivos concretos}
\label{sec:obxectivos}

Para acadar os obxectivos xerais expostos e acadar unha solución axeitada definíronse os seguintes obxectivos concretos:

\begin{itemize}
    \item Desenvolver un emulador do modelo orixinal de Game Boy, o DMG-01 (Dot Matrix Game).
    \item Testear os seus compoñentes fundamentais contra un conxunto de probas de referencia amplamente usado pola comunidade.
    \item Garantir o seu correcto funcionamento nos xogos máis populares da consola, usando unha mostra representativa de cartuchos con diferentes características de hardware.
    \item Deseñar e desenvolver unha interfaz axeitada para o seu uso en contornos embebidos.
    \item Compilar e optimizar o proxecto para asegurar a súa correcta execución en contornas de recursos limitados.
    \item Documentar unha construción de referencia sobre hardware accesible, que permita a calquera persoa montar, adaptar ou mellorar o seu propio dispositivo.
\end{itemize}

\section{Estructura da memoria}
\label{sec:estructura}
Esta memoria organízase nos seguintes capítulos:

\begin{enumerate}
    \item \textbf{Introdución}: contextualízase o proxecto, expóñense os seus obxectivos e detállase a estrutura.
    \item \textbf{Estado da arte}: expoñerase o estado da emulación de Game Boy e da súa comunidade de desenvolvedores,
    \item \textbf{Metodoloxía e planificación}: describirase a metodoloxía de traballo seguida, a planificación temporal do proxecto e a estimación dos seus custos.
    \item \textbf{Tecnoloxías e ferramentas}: presentaranse as tecnoloxías e ferramentas empregadas no desenvolvemento e xustificarase a súa elección.
    \item \textbf{Arquitectura da consola}: describirase a arquitectura da Game Boy, os seus compoñentes e o seu funcionamento en conxunto.
    \item \textbf{Desenvolvemento do emulador}: detallarase o proceso de desenvolvemento do emulador, as decisións de deseño tomadas e os retos afrontados.
    \item \textbf{Probas e validación}: describirase o proceso de testeo e validación do emulador, os resultados obtidos.
    \item \textbf{Conclusión e traballo futuro}: expoñeranse o aprendizaxe adquirido durante o desenvolvemento do proxecto e a súa relación co grao. Abordarase o liñas de traballo futuro para a continuación do desenvolvemento.
\end{enumerate}
